\chapter{Kết Luận}
\label{chap:conclusion}

\section{Tóm tắt đóng góp}

Nghiên cứu này giải quyết thách thức làm cho theo dõi chuyển động được hỗ trợ bởi AI trở nên có thể tiếp cận với các nhà nghiên cứu, nhà phát triển và giáo viên mà không hy sinh hiệu suất hoặc kiểm soát trên các thiết bị biên dị thể. Luận văn này trình bày một framework mã nguồn mở toàn diện thống nhất tiền xử lý dữ liệu, huấn luyện mô hình và triển khai biên vào một quy trình đơn thân thiện với người dùng. Các đóng góp chính của hệ thống trải dài đổi mới thuật toán, kỹ thuật phần mềm và xác thực thực nghiệm:

\subsection{Đóng góp kỹ thuật}
\begin{enumerate}
    \item \textbf{Giao diện tiền xử lý tương tác mới lạ}: Cơ chế phân đoạn cửa sổ thời gian kéo thả đại diện cho công cụ chú thích trực quan đầu tiên cho dữ liệu cảm biến chuyển động, loại bỏ sự lựa chọn truyền thống giữa việc coding logic phân đoạn phức tạp và chấp nhận các phương pháp tự động hộp đen. Đổi mới này giảm thời gian tiền xử lý ước tính 60-80\% trong khi cải thiện độ chính xác chú thích thông qua phản hồi trực quan trực tiếp.

    \item \textbf{Hỗ trợ đa thuật toán toàn diện}: Không giống như các nền tảng thương mại giới hạn ở mạng nơ-ron hoặc các thuật toán ML chuyên biệt, framework này tích hợp năm thuật toán (Random Forest, SVM, Neural Network, PyTorch MLP, PyTorch 1D-CNN) với các API nhất quán. Tính linh hoạt này cho phép người dùng chọn các thuật toán khớp với các ràng buộc triển khai của họ---Random Forest cho tạo mẫu nhanh (huấn luyện 0,15s), PyTorch MLP/CNN cho độ chính xác tối đa (99,13\%), hoặc SVM cho khả năng diễn giải lý thuyết.

    \item \textbf{Kiến trúc triển khai đa thiết bị nhận biết tối ưu hóa}: Thay vì gắn với một bo mạch duy nhất, tầng generator được tổ chức theo ba trục---loại mô hình, nền tảng đích và backend triển khai---cho phép xuất các gói C/C++/MicroPython cho Arduino-compatible boards, ARM Cortex-M, ESP-IDF, Zephyr và các backend runtime thay thế như TFLite Micro hoặc ONNX Runtime. Bốn chế độ tối ưu hóa (\texttt{accuracy}, \texttt{balanced}, \texttt{speed}, \texttt{power}) điều chỉnh độ chính xác số học và chiến lược bộ đệm theo ràng buộc tài nguyên của nền tảng đích.

    \item \textbf{Xử lý mất cân bằng lớp}: Bằng cách làm cho \texttt{class\_weight='balanced'} là một cấu hình mặc định thay vì tham số tùy chọn, framework hệ thống giải quyết một thách thức quan trọng nhưng thường bị bỏ qua trong theo dõi chuyển động thực tế. Thực nghiệm chứng minh các lớp thiểu số (Walking Downstairs: 38 mẫu, Walking Upstairs: 35 mẫu) đạt recall 0,971--1,000 và F1 $\geq$ 0,949, xác nhận hiệu quả của chiến lược cân bằng.

    \item \textbf{Đảm bảo tương đồng huấn luyện-triển khai}: Luận văn này xác định và giải quyết một thách thức ít được nghiên cứu trong edge ML: sự không khớp giữa pipeline trích xuất đặc trưng huấn luyện (Python) và suy luận (C++). Quy trình xác minh 12 điểm kiểm tra đảm bảo tính nhất quán bit-exact cho các đặc trưng miền thời gian, bao gồm sửa lỗi zero-padding, thống nhất công thức kurtosis/skewness, và reorder đặc trưng tự động tại thời điểm tạo mã. Phát hiện này nhấn mạnh rằng tính minh bạch không chỉ là triết lý mã nguồn mở mà là yêu cầu kỹ thuật thiết yếu cho triển khai edge ML đáng tin cậy \cite{paleyes2022challenges}.

    \item \textbf{Tăng cường dữ liệu nhận biết lớp}: Framework tích hợp hệ thống tăng cường dữ liệu với năm phương pháp biến đổi đặc thù IMU \cite{um2017data}, cùng với chiến lược nhận biết lớp mới lạ tự động phân biệt hoạt động tĩnh và động. Hoạt động tĩnh chỉ nhận micro-jitter ($\sigma = 0{,}01$) thay vì biến đổi đầy đủ, bảo toàn đặc tính phân biệt ``gần bất động'' thiết yếu cho phân loại chính xác. Thiết kế này giải quyết vấn đề class confusion được quan sát khi áp dụng tăng cường đồng nhất.

    \item \textbf{Ngưỡng tin cậy cho từ chối hoạt động không xác định}: Mã C++ đã tạo tích hợp cơ chế từ chối dự đoán dựa trên ngưỡng tin cậy, trích xuất trực tiếp từ đầu ra mô hình hiện có (softmax/tỷ lệ bỏ phiếu) mà không yêu cầu mô hình phát hiện bất thường riêng biệt \cite{hendrycks2017baseline}. Phương pháp này không tốn thêm bộ nhớ, không cần dữ liệu huấn luyện bổ sung, và là bước phòng thủ thiết yếu cho triển khai thực tế nơi đầu vào ngoài phân phối là phổ biến.
\end{enumerate}

\subsection{Xác Thực Thực Nghiệm}
Các thử nghiệm toàn diện sử dụng tập dữ liệu Nhận Dạng Hoạt Động Con Người năm lớp xác thực hiệu quả của framework:
\begin{itemize}
    \item \textbf{Độ Chính Xác Mô Hình}: 97,38--99,13\% trên năm thuật toán cho bài toán 5 lớp, 100\% cho bài toán 3 lớp cơ bản---xác nhận hiệu quả của 63 đặc trưng bất biến hướng
    \item \textbf{Hiệu Suất Theo Từng Lớp}: F1-scores 0,949--1,000 trên tất cả các lớp hoạt động, bao gồm các lớp thiểu số (Walking Downstairs, Walking Upstairs), chứng minh khả năng phân loại đáng tin cậy
    \item \textbf{Tính Khả Thi Thời Gian Thực}: Triển khai thành công trên nền tảng tham chiếu Seeed XIAO nRF52840 (ARM Cortex-M4F @ 64~MHz), với dữ liệu tham khảo cho thấy PyTorch 1D-CNN đạt 99,4\% độ chính xác trên thiết bị
    \item \textbf{Định Vị Cạnh Tranh}: Cung cấp tính minh bạch hoàn toàn, không phụ thuộc đám mây, và chi phí bằng không so với các nền tảng thương mại (Edge Impulse \$20--100/tháng, SensiML \$99--500/tháng)
    \item \textbf{Đa Thiết bị theo Kiến trúc}: Tầng generator hỗ trợ các đích Arduino-compatible, ARM Cortex-M, generic C/C++, ESP-IDF, MicroPython và Zephyr; Seeed XIAO được dùng như ca benchmark đại diện chứ không phải thiết bị đích duy nhất
    \item \textbf{Tương Đồng Huấn Luyện-Triển Khai}: Xác minh 12 điểm kiểm tra đảm bảo đặc trưng miền thời gian giống hệt giữa Python và C++, loại bỏ nguồn gốc chính của suy giảm hiệu suất triển khai biên
    \item \textbf{Tăng Cường Nhận Biết Lớp}: Chiến lược tăng cường phân biệt tĩnh/động ngăn chặn class confusion được quan sát với tăng cường đồng nhất, bảo toàn đặc tính phân biệt của hoạt động tĩnh
    \item \textbf{Từ Chối Hoạt Động Không Xác Định}: Ngưỡng tin cậy tích hợp cho phép thiết bị biên từ chối dự đoán khi không đủ tin cậy, tăng độ tin cậy cho triển khai thực tế
\end{itemize}

\subsection{Tác Động Rộng Hơn}
Ngoài các đóng góp kỹ thuật, công việc này thúc đẩy dân chủ hóa edge AI bằng cách:
\begin{itemize}
    \item \textbf{Loại Bỏ Rào Cản Kinh Tế}: Thay thế các đăng ký thương mại \$20-500/tháng bằng một lựa chọn thay thế mã nguồn mở miễn phí có thể tiếp cận với các nhà nghiên cứu trên toàn thế giới
    \item \textbf{Hạ Thấp Rào Cản Kỹ Thuật}: Cho phép người dùng không có chuyên môn hệ thống nhúng triển khai các hệ thống theo dõi chuyển động cấp chuyên nghiệp thông qua các quy trình làm việc GUI trực quan
    \item \textbf{Thúc Đẩy Tính Minh Bạch}: Cung cấp khả năng hiển thị đầy đủ vào tiền xử lý, trích xuất đặc trưng và tạo mã—quan trọng cho khả năng tái tạo khoa học và sử dụng giáo dục
    \item \textbf{Kích Thích Đổi Mới}: Kiến trúc có thể mở rộng cho phép các nhà nghiên cứu tích hợp các thuật toán, nền tảng và chiến lược tối ưu hóa mới lạ mà không bắt đầu từ đầu
\end{itemize}

\section{Xem Lại Các Mục Tiêu Nghiên Cứu}

Các mục tiêu nghiên cứu ban đầu được nêu trong Mục 1.3 đã được giải quyết toàn diện:

\begin{enumerate}
    \item \textbf{Hệ Thống Tiền Xử Lý Tương Tác} \checkmark: Giao diện cửa sổ thời gian kéo thả được triển khai thành công và tích hợp với kiểm tra tín hiệu trực quan

    \item \textbf{Hỗ Trợ Đa Thuật Toán} \checkmark: Năm thuật toán (RF, SVM, NN, PyTorch MLP, PyTorch 1D-CNN) được tích hợp, tất cả đạt 97,38--99,13\% trên bài toán 5 lớp, cho phép lựa chọn theo ràng buộc triển khai

    \item \textbf{Tạo Mã Không Phụ Thuộc Nền Tảng} \checkmark: Kiến trúc trình tạo modular được tổ chức theo ma trận model-platform-backend, hỗ trợ nhiều họ đích triển khai; benchmark định lượng hiện được báo cáo trên nền tảng tham chiếu ARM Cortex-M4 với bốn chế độ tối ưu hóa

    \item \textbf{Xác Thực Hiệu Quả Framework} \checkmark: Đánh giá toàn diện chứng minh độ chính xác cạnh tranh (97,38--99,13\%), triển khai thành công trên thiết bị tham chiếu (99,4\% trên thiết bị cho CNN bài toán 3 lớp), và khả năng tái đóng gói cùng một pipeline cho nhiều đích triển khai khác nhau

    \item \textbf{Khả Năng Tiếp Cận} \checkmark: GUI dựa trên web với đường cong học tập thấp làm cho phát triển edge AI tiên tiến có thể tiếp cận với phi chuyên gia trong khi duy trì khả năng cấp chuyên nghiệp

    \item \textbf{Độ Tin Cậy Triển Khai} \checkmark: Quy trình xác minh tương đồng huấn luyện-triển khai 12 điểm, ngưỡng tin cậy và tăng cường dữ liệu nhận biết lớp đảm bảo mô hình hoạt động đáng tin cậy trên phần cứng thực, không chỉ trong môi trường mô phỏng
\end{enumerate}

\section{Các Ứng Dụng Thực Tế}

Framework đã xác thực cho phép các ứng dụng thực tế đa dạng:

\paragraph{Y Tế và Phục Hồi Chức Năng}
Các hệ thống đeo để giám sát mức độ hoạt động của bệnh nhân, theo dõi tiến trình phục hồi chức năng (phục hồi di động sau phẫu thuật, trị liệu đột quỵ), phát hiện té ngã trong quần thể người cao tuổi và cung cấp dữ liệu hoạt động khách quan cho các thử nghiệm lâm sàng. Kiến trúc tiêu thụ thấp của vi điều khiển ARM Cortex-M4F hỗ trợ giám sát liên tục nhiều ngày mà không gây gánh nặng cho bệnh nhân.

\paragraph{Thể Thao và Thể Dục}
Phát hiện và phân loại tập luyện tự động (chạy, đạp xe, cử tạ), phân tích hình thức để phòng ngừa chấn thương (phát hiện bất đối xứng dáng đi, kỹ thuật nâng không đúng) và phản hồi đào tạo cá nhân hóa. Cửa sổ trượt 1,5 giây với bước trượt 0,75 giây đảm bảo phản hồi gần thời gian thực cho các ứng dụng huấn luyện.

\paragraph{Nghiên Cứu và Giáo Dục}
Nền tảng chi phí thấp để giảng dạy các khái niệm edge AI trong các khóa học đại học/sau đại học, cho phép các phòng thí nghiệm thực hành mà không cần giấy phép thương mại đắt đỏ. Các nhà nghiên cứu có thể nhanh chóng tạo mẫu các thuật toán nhận dạng hoạt động mới lạ và triển khai lên phần cứng để xác thực hiện trường. Bản chất mã nguồn mở tạo điều kiện cho nghiên cứu có thể tái tạo và đóng góp cộng đồng.

\paragraph{Nhà Thông Minh và Hỗ Trợ Sống Môi Trường}
Giám sát hoạt động cho người cao tuổi sống độc lập (phát hiện các mẫu không hoạt động bất thường, theo dõi thói quen hàng ngày), tự động hóa nhà thông minh nhận biết bối cảnh (điều chỉnh ánh sáng/nhiệt độ dựa trên hoạt động được phát hiện) và phân tích xu hướng sức khỏe dài hạn. Đặc tính tiêu thụ thấp của lớp vi điều khiển ARM Cortex-M hỗ trợ cảm biến luôn bật phù hợp cho các ứng dụng giám sát dài hạn.

\section{Các Hạn Chế và Công Việc Tương Lai}

Trong khi framework chứng minh khả năng mạnh mẽ, các hạn chế được thừa nhận vạch ra các hướng cho nghiên cứu tương lai:

\subsection{Sự Đa Dạng Tập Dữ Liệu}
Xác thực hiện tại sử dụng dữ liệu đối tượng đơn cho năm hoạt động. Công việc tương lai nên:
\begin{itemize}
    \item Tiến hành các nghiên cứu đa đối tượng trải dài các nhân khẩu học đa dạng (tuổi, giới tính, mức độ di động)
    \item Mở rộng phân loại hoạt động để bao gồm các sự kiện té ngã, các hoạt động phức tạp (nấu ăn, lái xe) và chuyển đổi hoạt động
    \item Xác thực trên các tập dữ liệu benchmark công khai (UCI HAR, WISDM, PAMAP2) cho các so sánh chuẩn hóa
\end{itemize}

\subsection{Các Mở Rộng Thuật Toán}
\begin{itemize}
    \item \textbf{Học Sâu Mở rộng}: Tích hợp LSTM/Transformer và mở rộng đường dẫn TFLite Micro cho các kiến trúc phức tạp hơn, hỗ trợ học từ đầu đến cuối từ dữ liệu cảm biến thô
    \item \textbf{Học Trực Tuyến}: Cho phép các mô hình đã triển khai thích ứng với các mẫu cụ thể của người dùng theo thời gian mà không cần kết nối đám mây
    \item \textbf{Các Phương Pháp Ensemble}: Kết hợp nhiều dự đoán thuật toán (bỏ phiếu RF + NN + SVM) cho độ bền vững được cải thiện
    \item \textbf{Học Không Giám Sát}: Khám phá hoạt động tự động mà không cần gán nhãn thủ công cho các ứng dụng khám phá
\end{itemize}

\subsection{Mở Rộng Nền Tảng}
\begin{itemize}
    \item Mở rộng ma trận benchmark liên thiết bị trên AVR (Arduino Uno), RISC-V (ESP32-C3), ARM Cortex-M0+ (Raspberry Pi Pico) và các board hỗ trợ MicroPython/Zephyr để chuyển hỗ trợ ở mức generator thành số liệu định lượng xuyên nền tảng
    \item Hỗ trợ các mục tiêu FPGA và ASIC cho các thiết bị đeo công suất siêu thấp (<1mW)
    \item Triển khai di động (Android, iOS) cho theo dõi chuyển động dựa trên điện thoại thông minh
\end{itemize}

\subsection{Các Tính Năng Nâng Cao}
\begin{itemize}
    \item \textbf{Hợp Nhất Cảm Biến}: Tích hợp đa phương thức (IMU + GPS + nhịp tim + âm thanh) với đồng bộ hóa đặc trưng tự động
    \item \textbf{AI Có Thể Giải Thích}: Trực quan hóa các mẫu cảm biến nào kích hoạt dự đoán cụ thể (tích hợp LIME, SHAP)
    \item \textbf{Học Liên Bang}: Huấn luyện mô hình cộng tác bảo vệ quyền riêng tư trên nhiều thiết bị
    \item \textbf{Backend Đám Mây}: Sao lưu dữ liệu tùy chọn, phiên bản mô hình và chia sẻ mô hình cộng đồng trong khi duy trì kiến trúc local-first
    \item \textbf{Hiệu Chỉnh Tin Cậy}: Áp dụng temperature scaling \cite{guo2017calibration} hoặc Platt scaling để cải thiện chất lượng hiệu chỉnh tin cậy softmax, và điều chỉnh ngưỡng $\tau$ tự động dựa trên đường cong ROC tập xác thực
    \item \textbf{Tăng Cường Thích Ứng}: Mở rộng chiến lược nhận biết lớp để tự động xác định tham số tăng cường tối ưu cho từng lớp dựa trên đặc tính tín hiệu (phương sai, entropy), thay vì chỉ phân biệt nhị phân tĩnh/động
\end{itemize}

\section{Suy Nghĩ Cuối Cùng}

Luận văn này đã trình bày một framework toàn diện giải quyết các thách thức cơ bản trong theo dõi chuyển động dựa trên biên: khả năng tiếp cận, hiệu suất và tính linh hoạt. Thông qua tiền xử lý tương tác mới lạ, xử lý mất cân bằng lớp hệ thống và tạo mã nhận biết tối ưu hóa, đã tạo ra một công cụ khớp với dễ sử dụng của các nền tảng thương mại trong khi cung cấp kiểm soát đầy đủ và chi phí bằng không. Giá trị cốt lõi của framework không nằm ở việc tối ưu cho riêng Seeed XIAO, mà ở khả năng chuyển cùng một pipeline sang nhiều thiết bị biên khác nhau mà vẫn giữ được tính minh bạch và khả năng kiểm chứng. Xác thực thực nghiệm xác nhận rằng dân chủ hóa và xuất sắc hiệu suất không loại trừ lẫn nhau—chúng có thể và nên cùng tồn tại.

Hành trình từ dữ liệu cảm biến thô đến mô hình edge AI đã triển khai liên quan đến nhiều bước, mỗi bước truyền thống yêu cầu chuyên môn chuyên biệt. Framework này thu gọn các rào cản đó, cho phép bất kỳ ai có kiến thức miền (các hoạt động cần phát hiện) tạo ra các hệ thống sẵn sàng triển khai mà không trở thành chuyên gia về xử lý tín hiệu, lý thuyết học máy hoặc lập trình hệ thống nhúng. Đây là bản chất của dân chủ hóa: hạ thấp rào cản mà không hạ thấp tiêu chuẩn.

Khi chúng ta nhìn về tương lai, sự tiến hóa tiếp tục của edge AI sẽ phụ thuộc vào các công cụ trao quyền thay vì hạn chế, giáo dục thay vì che giấu và bao gồm thay vì loại trừ. Framework này đại diện cho một bước trong hướng đó. Con đường phía trước dài, nhưng điểm đến—một thế giới nơi các khả năng AI tiên tiến có thể tiếp cận với tất cả mọi người—đáng để theo đuổi.

Mã nguồn, tài liệu và các tập dữ liệu ví dụ có sẵn công khai để hỗ trợ khả năng tái tạo, giáo dục và đóng góp cộng đồng. Các nhà nghiên cứu và nhà thực hành được khuyến khích xây dựng trên nền tảng này, mở rộng khả năng của framework và áp dụng nó cho các thách thức theo dõi chuyển động mới lạ. Cùng nhau, chúng ta có thể bắc cầu khoảng cách giữa nghiên cứu phòng thí nghiệm và triển khai thực tế, chuyển đổi bối cảnh phát triển edge AI.
